Theorem 1. Specifically, we remove noisy causal patterns
by following a greedy strategy, i.e., iteratively removing
the noisy causal pattern such that unfaith(wΩ) is minimized
in each step. In this way, we just use the set of retained
causal patterns, denoted by Ω, to approximate the output, i.e.,
v(x)≈Y (x) =∑
S∈ΩwS·CS(x) =∑
S∈ΩwS.
Ratio of the explained causal effectsRΩ. We propose a
metricRΩ to quantify the ratio of the explained salient causal
effects in Ω to the overall network output.
RΩ =
∑
S∈Ω|wS|∑
S∈Ω|wS| +|∆| (7)
where ∆=v(x)−∑
S∈ΩwS denotes effects of the unexplained
causal patterns.
Third, discovering that adversarial training boosts
the conciseness. As discussed in Section 4.3, we also dis-
cover that adversarial training [46] makes the DNN encode
3The setting ofτ is introduced in Section 4.2. Please see Appendix E
for more discussions
more sparse causal patterns than standard training, thus
boosting the conciseness of the causal graph.
3.3. Rewriting the causal graph as an AOG
The AOG is a hierarchical graphical model that encodes
how semantic patterns are formed for inference, which
has been widely used for interpretable knowledge repre-
sentation [42, 86], object detection [64], etc. In this sec-
tion, we show that the above causal graph can be rewrit-
ten into an And-Or graph (AOG), which summarizes com-
mon coalitions shared by different causal patterns to fur-
ther simplify the explanation. According to the SCM in
Eq. (2), the causal graph in Section 3.1 actually repre-
sents the And-Sum representation encoded by the DNN,
i.e.,v(x)≈∑
S∈ΩwS·CS(x)=∑
S∈ΩwS. In fact, such And-
Sum representation can be equivalently transformed into an
AOG.
The structure of a simple three-layer AOG is shown in
Fig. 1(c). Just like the causal graph in Fig. 1(b), at the bottom
layer of the AOG in Fig. 1(c), there aren leaf nodes represent-
ingn variables of the input sample. The second layer of the
AOG has multiple AND nodes, each representing the AND
relationship between its child nodes. For example, the AND
node x4x5x6 indicates the causal pattern S ={x4,x 5,x 6}
with the causal effect wS = 2.0. The root node is a noisy
OR node (as discussed in [42]), which sums up effects of
all its child AND nodes to mimic the network output, i.e.,
output =∑
S∈ΩwS·CS.
Furthermore, in order to simplify the AOG, we extract
common coalitions shared by different causal patterns as
new nodes to construct a deeper AOG. For example, in
Fig. 1(c), input variables x5 and x6 frequently co-appear
in different causal patterns. Thus, we consider x5,x 6 as a
coalition and add an AND node β ={x5,x 6} to represent
their co-appearance. Accordingly, the pattern {x4,x 5,x 6}
is simplified as{x4,β} (see Fig. 1(d)). Therefore, for each
coalition / causal patternS in an intermediate layer, its trig-
gering stateCS =∏
S′∈Child(S)CS′, where Child(S) denotes
all input variables or coalitions composingS. I.e., each coali-
tion / causal patternS is triggered if and only if all its child
nodes in Child(S) are triggered.
In order to extract common coalitions, we use the mini-
mum description length (MDL) principle [27] to learn the
AOGg as the simplest description of causal patterns. The
MDL is a classic way of summarizing patterns from data for
decades, which has solid foundations in information theory.
Given an AOGg and input variablesN , letM=N∪ Ωcoalition
denote the set of all leaf nodes and AND nodes in the bottom
two layers, e.g.M=N∪ Ωcoalition ={x1,x 2,...,x 6}∪{α,β} in
Fig. 1(d). The objective of minimizing the description length
L(g,M) is given as follows.
min
M
L(g,M) s.t. L(g,M) =L(M) +LM(g), (8)
5
